\section{Introduction}
\label{sec:intro}

Climate policy is a macrofinancial event. A fossil-fuel ban writes down capital that somebody owns and somebody else has lent against; a retrofit subsidy is a fiscal flow that lands on households with a particular distribution of housing stock; a decarbonisation path that misses its target does so through investment decisions taken under credit rationing. Most macroeconomic models cannot say any of this, because their accounting stops short of balance sheets or their environmental block is a reduced-form elasticity bolted onto a supply-side core. Ecological stock-flow consistent (E-SFC) models are built for exactly this joint question, and DEFINE-UK is the United Kingdom's.

DEFINE-UK is documented in the \emph{DEFINE-UK Model Manual, Version 1.1} \citep{georgedafermos2026manual}, released April 2026 alongside version~1.0. It is a closed monetary accounting system covering households, non-financial corporations, the power-generation sector, monetary and non-monetary financial institutions, government and the rest of the world, with a non-sectoral production module, coupled to an energy and emissions block, and it carries two UK-specific structures --- a housing stock tracked by energy efficiency and a power sector with separate fossil and non-fossil capital --- through which climate policy reaches the financial system rather than merely the output gap. Version~1.1 added regulatory policies that directly depreciate affected capital (stranded assets as an explicit balance-sheet event), limited forward-looking behaviour towards policy announcements, direct government ownership of non-fossil power capital, and a green sovereign bond.

Within PolicyEngine Macro \citep{policyengine}, an open suite of macroeconomic models for the UK and US, DEFINE-UK occupies a role no other member can fill: it is the only model in the suite that answers \emph{who bears a climate-policy scenario}. It is also the suite's most constrained member, and this paper is as much a record of those constraints as of what was reproduced.

\paragraph{What this paper is.} The suite's standard is that a model earns its place by replication: run the published code or reimplement the published equations, then compare against the published results at recorded tolerances, and report the comparison with numbers rather than adjectives --- the same bar applied to the Bank of England SVAR \citep{ahmadi2026boesvar} and the OBR macroeconometric emulator \citep{ahmadi2026obr}. This paper applies that standard to DEFINE-UK and records, target by target, where it passes, where it diverges, and where the publications themselves impose a ceiling on what can be checked at all.

\paragraph{Three constraints shape everything that follows.} The first is legal. The upstream repository is public but carries \textbf{no licence}, which means its code cannot be vendored, redistributed or hosted. The adapter therefore fetches it at pinned commit \texttt{846081a} at runtime and executes it as published, in R, unmodified --- and DEFINE-UK is the only member of the suite that is local-only, never hosted. The second is evidential. No machine-readable numeric scenario results are published for version~1.1: the manual's own results stop at the baseline table plus scenario \emph{design} parameters, and both companion papers are access-restricted. There is, in short, nothing to replicate on the scenario side, and Section~\ref{sec:validation} says so rather than manufacturing a comparison. The third is calibrational. The published baseline sits far from UK outturns --- 2025 real GDP growth of 4.66 per cent against an ONS 1.31 --- and the manual states plainly that its baseline ``should also not be seen as a prediction or forecast''. Taking that literally is what makes the adapter serve annualised scenario-minus-baseline \emph{deltas, never levels}.

\paragraph{Contribution.} The contribution is threefold, and none of it is a forecast.

First, \emph{a replication with its ceiling stated}. Section~\ref{sec:validation} reports four published-target gates. The manual's Table~4 macro block passes. Its emissions path diverges, and is pinned at the observed ratios so that further drift fails loudly rather than being absorbed. The two scenario targets are closed at the achievable ceiling and are explicitly not passes. Throughout, checks that cannot fail are separated from checks that can, at the point of claim: the 2025 growth pass clears a tolerance set after the run by 0.01\,pp, and the two conference-vintage anchors are bands we chose to fit a verbal quantity, one of which admits values that contradict the sentence it is named for.

Second, \emph{a calibration and a multiplier reported with their non-comparabilities attached}. Section~\ref{sec:calibration} gives the baseline against ONS, DESNZ and OBR figures as a computed, regression-tested artifact rather than prose, and reports a cumulative green-public-investment multiplier of 1.78 that is ours, not the authors', and that is not like-for-like against any published multiplier on horizon, denominator or closure. A previously quoted IMF comparator is recorded as withdrawn.

Third, \emph{a clean-room reimplementation in progress}. Because the upstream cannot be hosted, the route to a hostable DEFINE-UK is a from-scratch Python implementation of the published equations, with the upstream used strictly as a numerical output oracle and never read to write equations. Section~\ref{sec:reimplementation} reports milestone~1 passed and milestone~2 in progress --- 118 contiguous equations, Eqs.~(21)--(138) --- and treats the reimplementation as an instrument in its own right: reading the manual closely enough to run it has surfaced thirteen defects, which are carried as machine-readable gap records rather than papered over, and an audit of those records re-characterised four of them and withdrew one claim outright.

\paragraph{Road map.} Section~\ref{sec:literature} places the model in the SFC, ecological-macroeconomics and climate-finance literatures. Section~\ref{sec:model} sets out the model as the manual specifies it. Section~\ref{sec:source} identifies the source publications, the licence position, and the replication scope. Section~\ref{sec:implementation} describes the adapter and what it serves. Section~\ref{sec:validation} is the replication record. Section~\ref{sec:calibration} covers external calibration and the multiplier. Section~\ref{sec:reimplementation} covers the clean-room track. Section~\ref{sec:limitations} concludes with limitations. The appendix reproduces the published baseline table, the full gap-record inventory, and the dated run record.
